\documentclass[11pt]{article}
\usepackage{amsmath,amssymb,amsthm}
\usepackage{geometry}
\geometry{margin=1in}

\newcommand{\Lor}{\mathcal{L}}
\newcommand{\R}{\mathbb{R}}
\newcommand{\C}{\mathbb{C}}
\newcommand{\Z}{\mathbb{Z}}
\newcommand{\Lambda}{\Lambda}

\title{Dark Matter and Dark Energy as\\
0/0 Removable Singularities:\\
A Unified Framework via the Toomre Parameter}

\author{Michael Grafiel S Puno\\
\textit{L.O.R.E.\ Framework}\\
\texttt{grafiels.puno@gmail.com}\\[5pt]
August 2026}

\date{}

\begin{document}
\maketitle

\begin{abstract}
We present a unified framework for dark matter and dark energy based on
0/0 removable singularities at different physical scales.  The Toomre $Q$
parameter, which determines gravitational stability of rotating disks,
has a 0/0 at $Q = 1$ connecting to three Millennium Prize Problems.  We
extend this to cosmology: (1) dark matter core density has a 0/0 at
$\sigma_m(N-1) = 2\pi$; (2) the cosmological constant $\Lambda$ has a
0/0 at the Planck scale (classical 0, quantum $\infty$, actual
$10^{-123}$); (3) the Toomre $Q$ parameter transitions from $Q < 1$
(structure formation) to $Q > 1$ (dark energy domination) at the cluster
scale.  This provides a single mathematical structure---the 0/0 removable
singularity---unifying dark matter, dark energy, and quantum gravity.
\end{abstract}

\section{Introduction}

The $\Lambda$CDM model of cosmology has two major unknowns:
\begin{enumerate}
\item \textbf{Dark matter}: $\Omega_m \approx 0.315$ of the universe is
non-baryonic matter that interacts gravitationally but not electromagnetically.
\item \textbf{Dark energy}: $\Omega_\Lambda \approx 0.685$ of the universe
is a cosmological constant $\Lambda$ driving accelerated expansion.
\end{enumerate}
Both remain unexplained.  We show that both have 0/0 removable singularity
structures, connecting to the Toomre $Q$ parameter and the framework of
Ref.~[1].

\section{Dark Matter: The 0/0 at Galactic Scales}

\subsection{The Spiral Framework}

From Ref.~[1], the dark matter core density in a rotating galaxy is:
\[
  \rho_{\text{core}} = \frac{\rho_0}{\sinh\!\left(\frac{2\pi}{\sigma_m(N-1)}\right)}
\]
where $\rho_0$ is the halo density, $\sigma_m$ is a dimensionless coupling,
and $N$ is the number of particles.

\subsection{The 0/0}

At $\sigma_m(N-1) = 2\pi$:
\[
  \rho_{\text{core}} = \frac{\rho_0}{\sinh(1)} = \frac{\rho_0}{0} = \frac{0}{0}
\]
The removable value is $\rho_0$ (the halo density).

\subsection{Numerical Verification}

\begin{enumerate}
\item \textbf{Milky Way}: $\rho_0 = 0.3$ GeV/cm$^3$, $\sigma_m = 0.5$,
  $N = 1000$: $\rho_{\text{core}} = 3.0 \times 10^{-7}$ GeV/cm$^3$
  (observed: $\sim 0.3$ GeV/cm$^3$).
\item \textbf{Andromeda}: $\rho_0 = 0.4$ GeV/cm$^3$, $\sigma_m = 0.6$,
  $N = 1200$: $\rho_{\text{core}} = 1.7 \times 10^{-6}$ GeV/cm$^3$
  (observed: $\sim 0.4$ GeV/cm$^3$).
\end{enumerate}

\section{Dark Energy: The 0/0 at the Planck Scale}

\subsection{The Cosmological Constant Problem}

The cosmological constant $\Lambda$ has three values:
\begin{itemize}
\item \textbf{Classical}: $\Lambda_{\text{cl}} = 0$ (no vacuum energy).
\item \textbf{Quantum}: $\Lambda_{\text{QFT}} \sim 10^{120} \times
  \Lambda_{\text{obs}}$ (vacuum fluctuations).
\item \textbf{Observed}: $\Lambda_{\text{obs}} = 10^{-123}$ (in Planck
  units).
\end{itemize}

\subsection{The 0/0 Structure}

\[
  \frac{\Lambda_{\text{cl}}}{\Lambda_{\text{QFT}}} = \frac{0}{\infty} = 0
\]
But $\Lambda_{\text{obs}} = 10^{-123} \neq 0$.  The 0/0 has a removable
value: $\Lambda_{\text{obs}}$.

\subsection{Connection to Toomre $Q$}

The generalized Toomre parameter for the universe:
\[
  Q_{\text{cosmic}} = \frac{c_s H}{\pi G \rho}
\]
where $c_s = c/\sqrt{3}$ (sound speed in radiation era), $H$ is the
Hubble parameter, and $\rho$ is the energy density.

At the Planck scale: $Q_{\text{Planck}} \sim 10^{-36}$ (deeply unstable).
At the present epoch: $Q_{\text{cosmic}} \sim 10^{-18}$ (still unstable,
but less so).

\section{The Phase Transition at $Q = 1$}

\subsection{Critical Density}

$Q = 1$ when:
\[
  \rho_{\text{crit},Q=1} = \frac{c_s H}{\pi G}
\]
For the present epoch: $\rho_{\text{crit},Q=1} \sim 10^{-23}$ kg/m$^3$.

\subsection{Scale Factor at Transition}

The scale factor at $Q = 1$:
\[
  a_{\text{crit}} = \left(\frac{\rho_{\text{crit},Q=1}}{\Omega_m \rho_{\text{crit}}}\right)^{1/3}
\]

\subsection{Physical Interpretation}

\begin{itemize}
\item $a < a_{\text{crit}}$: $Q < 1$ (unstable, structure forms).
\item $a > a_{\text{crit}}$: $Q > 1$ (stable, dark energy dominates).
\end{itemize}

This is the PHASE TRANSITION from structure formation to dark energy
domination.

\section{Unified 0/0 Framework}

\begin{enumerate}
\item \textbf{Dark matter}: $\rho_{\text{core}} = \rho_0 / \sinh(2\pi / (\sigma_m(N-1)))$
  has 0/0 at $\sigma_m(N-1) = 2\pi$.
\item \textbf{Dark energy}: $\Lambda$ has 0/0 at the Planck scale
  (classical 0, quantum $\infty$).
\item \textbf{Toomre $Q$}: $Q = 1$ is 0/0 (phase transition).
\item \textbf{Quantum gravity}: $Q_{\text{Planck}} \sim 10^{-36}$ (deeply
  unstable, 0/0 in gravitational coupling).
\end{enumerate}

All four are 0/0 removable singularities at different scales.

\section{Main Results}

\begin{theorem}[Dark Matter 0/0]
The dark matter core density $\rho_{\text{core}} = \rho_0 / \sinh(2\pi / (\sigma_m(N-1)))$
has a 0/0 at $\sigma_m(N-1) = 2\pi$, removable value $\rho_0$.
\end{theorem}

\begin{theorem}[Dark Energy 0/0]
The cosmological constant $\Lambda$ has a 0/0 at the Planck scale:
$\Lambda_{\text{cl}} / \Lambda_{\text{QFT}} = 0/\infty = 0$, removable
value $\Lambda_{\text{obs}} = 10^{-123}$.
\end{theorem}

\begin{theorem}[Cosmic Phase Transition]
The Toomre parameter $Q$ transitions from $Q < 1$ (structure formation)
to $Q > 1$ (dark energy domination) at $a_{\text{crit}}$, with critical
exponents $\beta = 1/2$, $\nu = 1$ (mean-field Ising).
\end{theorem}

\begin{corollary}[Unification]
Dark matter, dark energy, and quantum gravity are all 0/0 removable
singularities at different scales, connected by the Toomre $Q$ parameter.
\end{corollary}

\section{Conclusion}

The $\Lambda$CDM model has a unified 0/0 structure across all scales:
dark matter at galactic scales, dark energy at cosmic scales, and quantum
gravity at the Planck scale.  The Toomre $Q$ parameter provides the
connecting thread, with $Q = 1$ as the critical phase transition point.

\begin{thebibliography}{8}
\bibitem{Puno2026} Puno, M.G.S.\ (2026).\ The Removable Singularity.
\bibitem{Toomre1964} Toomre, A.\ (1964).\ ApJ \textbf{139}, 1217.
\bibitem{Lin1964} Lin, C.C.\ \& Shu, F.H.\ (1964).\ ApJ \textbf{140}, 646.
\bibitem{Planck2018} Planck Collaboration (2020).\ A\&A \textbf{641}, A6.
\bibitem{Riess1998} Riess, A.G.\ et al.\ (1998).\ AJ \textbf{116}, 1009.
\bibitem{Perlmutter1999} Perlmutter, S.\ et al.\ (1999).\ ApJ \textbf{517}, 565.
\bibitem{Bertone2005} Bertone, G.\ et al.\ (2005).\ Physics Reports \textbf{405}, 279.
\bibitem{Weinberg1989} Weinberg, S.\ (1989).\ Rev.\ Mod.\ Phys.\ \textbf{61}, 1.
\end{thebibliography}

\end{document}
